% Part: second-order-logic
% Chapter: syntax-and-semantics
% Section: semantic-notions

\documentclass[../../../include/open-logic-section]{subfiles}

\begin{document}

\olfileid[pt]{sol}{syn}{sem}
\olsection{Noções semânticas}

\begin{explain}
As noções lógicas centrais de \emph{validade}, \emph{consequência} e
\emph{satisfatibilidade} são definidas da mesma maneira para a lógica de
segunda ordem e para a lógica de primeira ordem, exceto pelo fato de que a
relação de satisfação subjacente agora é a das !!{formula}s de segunda
ordem. Uma !!{sentence} de segunda ordem é, naturalmente, !!a{formula} na
qual todas as variáveis, inclusive as variáveis de predicado e de função,
estão ligadas.
\end{explain}

\begin{defn}[Validade]
Uma sentença $!A$ é \emph{válida}, $\Entails !A$, sse $\Sat{M}{!A}$ para
toda !!{structure}~$\Struct M$.
\end{defn}

\begin{defn}[Consequência]
Um conjunto de sentenças~$\Gamma$ \emph{acarreta} uma sentença~$!A$,
$\Gamma \Entails !A$, sse, para toda !!{structure}~$\Struct M$ tal que
$\Sat{M}{\Gamma}$, temos $\Sat{M}{!A}$.
\end{defn}

\begin{defn}[Satisfatibilidade]
Um conjunto de sentenças~$\Gamma$ é \emph{satisfatível} se
$\Sat{M}{\Gamma}$ para alguma !!{structure}~$\Struct M$. Se $\Gamma$ não é
satisfatível, é chamado \emph{insatisfatível}.
\end{defn}

\end{document}

